%``The main tactic that cancerous tissue applies to evade the immune system, is interaction with it,''
%Mike van Santvoort writes in his PhD thesis.
%Normally, the immune system detects and eliminates misbehaving cells.

At the time of writing, cancer is still one of the leading causes of death in the Netherlands~\cite{noauthor_ranglijsten_nodate}.
Roughly half of Dutch citizens are diagnosed with cancer at some point in their lifetime~\cite{noauthor_kans_nodate}.
The human immune system is naturally capable of tracking, targeting and killing tumors.
Sadly, tumor cells can develop methods to defend themselves from immune cells.
Strategies include deactivating immune cells, evading detection, and keeping immune cells at a distance.
As a result of these ``defenses'', many therapies designed to eradicate cancer, including immunotherapy,
are ineffective for many tumors.

One of the ways that cancerous tissue tries to evade the immune system is through the PD-L1 protein.
PD-L1 is an important immune checkpoint, a brake that prevents our immune system from becoming overactive and attacking
healthy tissue.
Cancer cells can abuse PD-L1.
As a result, PD-L1 has become a target for new cancer treatments, with
FDA-approved immunotherapy targeting this protein~\cite{twomey_cancer_2021}.

For the PD-L1 protein to be synthesized, the first step is transcription of its regulating gene CD274.
To start transcription, certain proteins need to bind to the gene sequence.
These proteins are called transcription factors, or TF for short.
It is possible for a single TF to start transcription.
In other cases, a combination of TFs, often in the form of a ``complex'', is required.
If any of the TF in the complex is absent, transcription will not start.
A gene is typically transcribed by several TFs and/or combinations of TFs.
This creates an enormous search space that is difficult to traverse.

Still, biologists and clinicians are interested in when and how combinations of TFs start synthesis of proteins like PD-L1.
Who are the ``partners in crime'', helping PD-L1 be synthesized and ultimately helping cancer spread?
Of course, the TF that enable production of PD-L1 in cancer cells may be the same ones acting in healthy cells.
Still, for the sake of this investigation, we shall consider TFs acting together as ``partners in crime''
and embark on an investigative journey.
The knowledge gained could support identification of potential clinical targets, such that PD-L1 synthesis
may be prevented from the very beginning.

The aim of this research is to identify, from data, candidate groups of TF that may act together to enable PD-L1 production.
The candidate groups could then be studied further, for instance in the lab.
Understanding which groups are enabling the PD-L1 evasion tactic may ultimately help to develop strategies to prevent it.
For even a moderate number of TFs and a moderate group size, the number of possible groups is large.
Consequently, it is attractive or even necessary to study candidate groups systematically and computationally.

We wish to design computational methods capable of identifying these groups from bulk transcriptomics,
as such data is widely available.
We wish to use TF activity data for our analysis,
as a TF needs to be active to bind to a gene and contribute to transcription.
Measuring how active a TF is inside a cell, however, is difficult.
Readily available are bulk RNA sequencing data.
These measure the average number of RNA transcripts across an entire group of cells in a tissue sample.
The RNA transcript count will be simply referred to as ``gene expression''.
Gene expression is not equivalent to protein activity.
If gene expression is high, it does not necessarily mean that the resulting protein is active.
Intermediate steps can affect the protein activity, like protein synthesis and post-translational modifications.
Therefore, an existing method is used to estimate TF activity based on the gene expression profile in a tissue sample~\cite{badia-i-mompel_decoupler_2022}.

The activation of the CD274 gene is represented by the PD-L1 gene expression.
We refer this as simply ``PD-L1 expression'' and it will be the quantity of interest.
In biological modeling, it is common to assume an ``OR'' relation between TFs and target gene expression, i.e.,
``if TF A or TF B or TF C is active, then gene X is likely expressed''.
When investigating which TF act together inside a complex, the natural analogue is to assume an ``AND'' gate.
That is, ``if TF A and TF B are active, then gene X is likely expression''.
In reality, different complexes may lead to a gene's transcription (expression), so a more fitting model is an
``OR-of-ANDs''.
Denote an ``OR'' by $\vee$ and an ``AND'' by $\wedge.$
An example is
\begin{gather*}
    (\text{TF A is active} \wedge \text{TF B is active}) \lor (\text{TF C is active} \wedge \text{TF D is active})\\
    \Longrightarrow \text{Gene X is expressed}.\\
\end{gather*}
From a mathematical perspective, this leads us to the subject of ``ruleset learning''~\cite{yang_learning_2022}.
Ruleset learning is concerned precisely with learning a formula of this form from data.
A major advantage of rulesets over, for instance, neural networks, is that rulesets are easily interpretable.

There exist related data mining techniques known for their interpretability, like decision trees.
Rulesets and decision trees are mathematically related; one can always be expressed in the form of the other.
However, they differ conceptually and in terms of expressiveness.
A relatively simple ruleset can be equivalent to a large, redundant decision tree and vice versa.
Notably, decision trees hold a certain hierarchy; conditions checked higher up the tree are ``more important''.
The root node, for instance, plays a pivotal role in all decisions.
However, we are looking for a diverse set of relevant conditions.
It is not obvious that activity of one TF individually should govern the entire decision process.
Therefore, while substantial literature exists on optimal decision trees (\cite{van_der_linden_optimal_2026}),
we do not consider decision trees suitable for this problem.
Another related technique is optimal rule list mining, which was studied in~\cite{angelino_learning_2018}.
A rule list (or decision list) is an ordered list of if-statements, evaluated sequentially to make a prediction.
For instance,
\begin{gather*}
    \textbf{if } (\text{TF A is active} \wedge \text{TF B is active}) \textbf{ then } \text{predict ``Gene X expressed``},\\
    \textbf{else if } (\text{TF C is active}) \textbf{ then } \text{predict ``Gene X expressed``}
\end{gather*}
Again there is a hierarchy; later rules assume that earlier rules were false.

Therefore, we build upon earlier research into optimal rule sets.
In particular, we employ the techniques proposed in~\cite{dash_boolean_2020}
and later studied further in~\cite{lawless_interpretable_2023}.
We will study whether rule sets can be effective in predicting PD-L1 expression from TF activity scores
in the context of bulk transcriptomics.
Furthermore, we will investigate whether the rule sets describe an intrinsic phenomenon in the data.
Precisely, we study whether the rule sets found are consistent and generalize to unseen data.


\section{Data}\label{sec:data}
The gene expression datasets used are Broad and Sanger.
These were extracted from the ``RNA-Seq 2022-06-24'' data collection,
downloaded from the Cell Model Passports hub~\cite{van_der_meer_cell_2019, noauthor_cell_nodate}.
This collection contains RNA transcript counts of thousands of genes, including PD-L1, for hundreds of samples.
The measurements of each sample were done on cancer cell lines.
These are derived from human tissue and are often immortalized.
The advantage is that cancer cell lines behave more consistently in experiments compared to fresh human cancer tissue.
We refer to a sample as a ``datapoint'' or simply ``point''.
Combined, the datasets contain $1291$ points.
These are too few points to use methods likes neural networks.

For each sample, the tissue status and tissue type is also recorded in the dataset.
The samples consist of a wide variety of cancer types found in various types of tissue.
Tissue status among the samples is not homogenous.
Sample counts per dataset and tissue status are given in \autoref{tab:data_counts}.

\begin{table}[h]
    \centering
    \resizebox{\textwidth}{!}{%
\begin{tabular}{l rrrrr r}
\toprule
\textbf{Dataset} & \textbf{Tumor} & \textbf{Metastasis} & \textbf{Papiloma} & \textbf{Normal} & \textbf{Unknown} & \textbf{Total} \\
\midrule
\textbf{Broad}  & 507 & 272 & 1 & 1 & 68  & 849  \\
\textbf{Sanger} & 240 & 113 & 0 & 6 & 83  & 442  \\
\midrule
\textbf{Total}  & 747 & 385 & 1 & 7 & 151 & 1291 \\
\bottomrule
\end{tabular}
}
    \caption{Sample counts by dataset and tissue type.}
    \label{tab:data_counts}
\end{table}

Out of all candidate TF, a total of 18 were selected to be of interest.
This selection was based on literature, in particular two criteria:
(1) the TF can physically bind to the CD274 gene sequence and (2)
there is a recorded functional response in literature.
A recorded functional response means that some significant change in PD-L1 expression is associated
with the activity of the TF\@.

The TF activities are estimated using DecoupleR model~\cite{badia-i-mompel_decoupler_2022, badia-i-mompel_transcription_nodate},
using CollecTRI as prior knowledge~\cite{muller-dott_expanding_2023}.
DecoupleR estimates TF activity in a sample based on the full gene expression profile (thousands of genes).
Simply put, it estimates TF activity based on the expression (level of under- or overregulation)
of the genes it regulates (according to the prior knowledge).
The output is, for each TF, a ``TF activity score'', also referred to as ``TF activity'' or ``TF activity estimate''.
Descriptive statistics for the corresponding TF activities and PD-L1 expression in both Broad and Sanger
can be found in \autoref{tab:descr_broad} and \autoref{tab:descr_sanger}, respectively.



The cross-correlation matrix of TF activity scores is depicted in \autoref{fig:corr_activity}.
The cross-correlation matrix of the TF gene expressions (used to calculate TF activities) is depicted in \autoref{fig:corr_TPM}.
Notably, TF activity scores are highly correlated.
This indicates that the selected TFs tend to be active or inactive simultaneously.
This may be expected, if these TFs are known to regulate similar genes.
Meanwhile, gene expressions are not highly correlated.
This suggests that gene expression is not an accurate indicator of simultaneous TF activity.
Alternatively, TF activity scores could be biased towards showing correlation due to the method of estimation.
Both figures also include correlation with PD-L1 expression in the bottom row.
The PD-L1 expression shows relatively low correlation with TF activity scores and gene expressions, respectively.
This indicates that activity or expression of a single TF does not explain PD-L1 expression substantially.
This motivates us further to consider the effect of groups of TFs.


PD-L1 expression for Broad and Sanger dataset combined was normalized to mean $0$ and standard deviation $1$,
The histogram of PD-L1 expression of Broad samples is depicted in \autoref{fig:pdl1_expression_broad}.
It shows that the majority of observations clusters around the minimum.
Note that the plot starts at the minimum observed expression value and is capped at an expression value of $0.5$,
while the maximum observed PD-L1 expression is $3.55$.

\begin{figure}
    \centering
    \includesvg[width=\textwidth]{gen/data/pdl1_expression_broad.svg}
    \caption{Histogram of PD-L1 expression in Broad dataset. Increasing quantiles are also depicted from left to right using vertical lines.}
    \label{fig:pdl1_expression_broad}
\end{figure}

The goal is to infer general combinations of TF contributing to PD-L1 expression.
Hence, it is important that the rulesets produced generalize to other datasets.
This means rulesets need to be tested on unseen data.
Broad serves as the training set.
Cross-validation will be used on splits of Broad to study how well the results various methods, under various parameters, generalize.
Sanger serves as the holdout dataset and will be used only once, at the very end (\autoref{ch:holdout}).
This is the ultimate test whether the method is able to produce general rulesets.

%
%\section{Transcription Factor activity}\label{sec:data_TF_activity}
%Gene expression is not equivalent to protein activity.
%If gene expression is high, it does not necessarily mean that the resulting protein is active.
%There are intermediate steps that can affect the protein activity.
%For a TF to bind to a gene, however, the TF needs to be active.
%Therefore, using gene expression directly may not be suitable for our analysis.
%To address this issue, an estimate of TF activities is calculated using
%the decoupleR model with CollecTRI as prior knowledge network
%\cite{badia-i-mompel_decoupler_2022, muller-dott_expanding_2023, badia-i-mompel_transcription_nodate}.
%
%DecoupleR estimates TF activity based on the expression of all genes in the dataset via a collection of linear models. \todo{How much should I go into this model?}
%
%
%

%
%\subsection{Tissue status}\label{subsec:tissue_status}


\section{Outline}\label{sec:outline}
In \autoref{ch:boolean-model} we apply the method studied in~\cite{lawless_interpretable_2023} to
the Broad dataset to study whether high PD-L1 expression can be captured in a ruleset.
In \autoref{ch:threshold-model} we describe a version of the model tailored to thresholded continuous
features.
We use it to conduct additional experiments to understand the behavior of the model on the Broad dataset.
In \autoref{ch:ramp_model} we build upon the model of \autoref{ch:threshold-model} to study
whether an alternative loss function, the ramp loss, can help produce more consistent rulesets.
Based on these results, we perform final experiments in \autoref{ch:experiments}.
In particular, we consider subsets of the data, a different dataset and an experiment in which TF are excluded
to identify ``important'' TF\@.
The rulesets that perform adequately on the Broad dataset are tested on the holdout dataset,
Sanger, in \autoref{ch:holdout}.
The performance of these rulesets is presented alongside the best performing rulesets on the Sanger dataset.
We propose three recommendations for future research in \autoref{ch:recommendations}.
Finally, concluding remarks are covered in \autoref{ch:conclusion}.
